\subsection*{从电场力做功到电势}

静电力是保守力。把试探电荷 $q$ 从 $A$ 移到 $B$，电场力做功只由始末位置决定，定义
\[
  W_{A\to B}=q(\varphi_A-\varphi_B),
  \qquad U=q\varphi.
\]
电势是单位正电荷的势能，是描述电场的标量函数；静电势能则属于具体电荷与外电场组成的系统。
同一点的电势与放入多大试探电荷无关，势能却正比于试探电荷。若试探电荷为负，它受力方向与
电场相反，沿电场方向移动时势能反而增加。因此判断能量变化时不能只背“沿电场电势降低”，还要
同时考虑电荷符号。

电势零点可以任意选择，真正可测的是电势差。孤立有限电荷分布常取无穷远电势为零；无限长线
电荷或无限大带电平面的绝对电势会发散，只能选择有限参考点讨论电势差。多个点电荷产生的电势
直接作代数叠加：
\[
  \varphi(\vect r)=\frac1{4\pi\varepsilon_0}
  \sum_i\frac{q_i}{|\vect r-\vect r_i|}.
\]
因为电势是标量，复杂电荷分布往往先求电势再求梯度，比直接叠加电场矢量更简洁。

电场与电势的局部关系为
\[
  \vect E=-\nabla\varphi.
\]
负号表示电场指向电势下降最快的方向。直角坐标中三个分量分别是
$E_x=-\partial\varphi/\partial x$、$E_y=-\partial\varphi/\partial y$、
$E_z=-\partial\varphi/\partial z$。等势面上任意切向位移满足 $\dd\varphi=0$，所以电场只能
沿法向。图30.1把匀强电场、点电荷和偶极场并列，是为了同时展示平行面、同心球面和复杂曲面
三种几何，而不是只记住“等势线互相平行”这一种特殊情况。

\subsection*{典型分布的电势与等势面几何}

球对称电荷分布在球外的电势等效于总电荷集中在球心。均匀带电球壳内部电场为零，因此内部电势
处处等于球面电势，并不等于零。均匀带电实心球内部电场随 $r$ 线性变化，积分后电势是 $r$ 的
二次函数。计算时应先保证电势在分界面连续，再用无穷远零点确定常量。

两个异号点电荷的零电势面满足到两电荷距离之比为常量，是 Apollonius 球面。当 $+Q$ 位于
$x=a$、$-q$ 位于原点且 $Q>q$ 时，球心在原点左侧，球面包围较小的负电荷；不可能把整个球画在
两电荷之间。若 $Q=q$，距离比变为1，球面退化为中垂平面。图形必须与配方所得球心和半径一致。

两根异号无限长线电荷的横截面问题同样给出 Apollonius 圆。等势圆应包围对应的线电荷，电势
越接近线电荷，圆越小且越靠近该电荷。电场线则是通过两根导线的一族圆弧，并与等势圆正交。
如果画出的等势圆既没有包围正线电荷也没有包围负线电荷，即使形式上有许多圆，位置关系仍是
错误的。

有限长带电线段的电势可由
\[
  \varphi=k\lambda\int_{-a}^{a}
  \frac{\dd x'}{\sqrt{(x-x')^2+\rho^2}}
\]
求得。结果只依赖场点到两端点的距离组合，等势面为绕线段轴旋转的椭球面，配套的正交曲面为
双曲面。这种焦点几何与引力轨道中的焦点概念不同，但都要求图中的端点、焦点和距离定义严格
一致。

\subsection*{电子伏特、原子电离与气体击穿}

电子伏特把电势差和微观粒子能量直接连接起来：一个元电荷通过1伏电势差，能量改变1电子伏特。
原子的第一电离能通常为数电子伏特到二十几电子伏特，并随原子序数呈明显周期性。稀有气体电子
壳层较稳定，第一电离能形成峰值；碱金属最外层电子较易移走，形成低谷。图30.3(a)保留数值纵轴和
元素序列，才能让读者看到“多数处于 $5$--$25\,\mathrm{eV}$”并非随意估计。

核过程的典型能量进入 $\mathrm{MeV}$ 范围，高能加速器则使用 $\mathrm{GeV}$、$\mathrm{TeV}$。
加速器中的电场负责改变粒子能量，磁场主要负责偏转和聚焦。粒子在环形轨道中反复通过加速腔，
每次获得 $q\Delta\varphi$，许多次累积后达到高能；不能把环形隧道本身误解为始终沿切向做功的
静电场。

气体中电子在两次碰撞之间的平均自由程为 $l$，电场在这段距离内提供的能量约为 $eEl$。当它
超过原子激发能时可能发光，超过电离能时会产生新的电子和离子，电子数在连续碰撞中倍增，最终
形成击穿。闪电的明亮分支是电离后的导电通道，不是预先存在的电场线。压强、气体成分、自由程
和电极曲率都会改变击穿条件，所以 $eEl\sim10\,\mathrm{eV}$ 只能作为数量级判断。

\subsection*{电荷系统能量与电场能量}

把多个点电荷从无穷远缓慢搬运到给定位置，外力做功等于系统静电势能：
\[
  U=\sum_{i<j}\frac{q_iq_j}{4\pi\varepsilon_0r_{ij}}
   =\frac12\sum_i q_i\varphi_i^{(\mathrm{others})}.
\]
因子 $1/2$ 用来避免每一对相互作用被计算两次。连续分布形式
$U=\tfrac12\int\rho\varphi\,\dd V$ 也有同样含义。若把自能发散的理想点电荷包含在积分中，
还需明确是否剔除每个电荷与自身的相互作用。

电场本身具有能量，真空中的能量密度为
\[
  u_E=\frac12\varepsilon_0E^2,
  \qquad U=\int u_E\,\dd V.
\]
对均匀带电球壳，既可以用组装电荷的方法计算，也可以把球外电场能量积分；两种结果相同。后者
显示能量储存在有电场的空间，而不是只“附着”在电荷上。由能量对几何参数求导可以得到静电力，
但必须先说明控制条件：孤立系统通常保持电荷不变，接电源系统可能保持电势不变，两种条件下的
能量函数和外力功并不相同。

\subsection*{从电场积分到电势分布}

选定零势面以后，任意点电势由
\[
 \varphi(B)-\varphi(A)=-\int_A^B\vect E\cdot\dd\vect l
\]
确定。静电场中积分与路径无关，所以可以选择计算最方便的路线；若沿等势面移动，线元始终与
电场垂直，积分为零。反过来，等势面上的电势不变意味着电场没有切向分量，故电场线必与等势面
正交。画图时，场线与等势线若以斜角相交，便不是风格问题，而是违反
$\vect E=-\nabla\varphi$ 的物理错误。

对点电荷系，取无穷远为零可写成
\[
 \varphi(\vect r)=\frac1{4\pi\varepsilon_0}\sum_i
 \frac{q_i}{|\vect r-\vect r_i|}.
\]
电势是标量，叠加时只作代数加法；电场是矢量，需要逐分量相加。总电势为零的点不一定电场为零，
电场为零的点也不一定电势为零。前者表示各电荷的标量贡献抵消，后者表示电势在该点梯度为零，
两者条件完全不同。

无限长线电荷的绝对电势不能把无穷远作为零点，因为积分呈对数发散；只能规定某个有限参考半径
$r_0$ 的电势，写成
\[
 \varphi(r)-\varphi(r_0)=-\frac{\lambda}{2\pi\varepsilon_0}
 \ln\frac r{r_0}.
\]
两根线电荷组成的等势线由两个距离之比决定，几何上是 Apollonius 圆。每一簇闭合等势圆必须围住
对应的线电荷；若圆心、半径画错，使电荷落在圆外，便与距离比的几何定义矛盾。

\subsection*{势能、外力功与电场力功}

电荷 $q$ 在外电势 $\varphi$ 中的势能为 $U=q\varphi$。它从 $A$ 移到 $B$ 时，电场力做功
\[
 W_{\rm e}=q[\varphi(A)-\varphi(B)]=-\Delta U.
\]
若缓慢移动、始末动能不变，外力做功 $W_{\rm ext}=\Delta U$。若粒子自由加速，则
$\Delta K=-\Delta U$。这三种说法对应不同的功，符号不能混用。对电子 $q=-e$，它沿电场方向
移动时电势降低但势能升高；“正电荷趋向低电势”不能直接照搬给负电荷。

由势能曲线判断平衡时，$F_x=-\dd U/\dd x$。平衡点满足 $U'(x_0)=0$；若
$U''(x_0)>0$ 是稳定平衡，若 $U''(x_0)<0$ 是不稳定平衡。对受约束运动，只能沿允许方向取导数；
把三维空间的极值结论未经检查地套到一条轨道上，可能漏掉横向不稳定性。

\subsection*{微观能量尺度的读图方法}

第一电离能曲线不是要求逐个死记数值，而是把电子壳层结构转化为可比较的能量尺度。同一周期中，
核电荷增加总体使电子束缚增强，但亚层填充和电子配对会造成局部起伏；到稀有气体形成峰值，进入
下一周期的碱金属后骤降。图中的横轴必须按元素顺序连续排列，纵轴必须保留电子伏特量级，才能
同时读出周期性和数量级。

把电离能与气体放电联系起来时，还需区分电子一次碰撞获得的能量与外加电压。电压很高并不保证
单次自由程内获得足够能量；压强升高会缩短自由程，电场不均匀又使尖端附近更容易先击穿。因而
宏观电压、几何尺度和微观电离能必须通过 $eEl$ 联系，不能只比较两个表面数值。

电势零点的选择不影响电场，也不影响同一电荷在两点间的势能差；但讨论“某点电势为正还是为负”
时必须先说明参考零点。对有限局域电荷系通常取无穷远为零，对无限线电荷或无限平面则应选择有限
参考位置，只比较电势差。

\subsection*{连续带电体的电势积分}

连续电荷分布的电势由点电荷结果叠加得到：
\[
 \varphi(\vect r)=\frac1{4\pi\varepsilon_0}
 \int\frac{\rho(\vect r')}{|\vect r-\vect r'|}\,\dd V',
\]
线分布和面分布则把 $\rho\dd V'$ 换成 $\lambda\dd l'$、$\sigma\dd S'$。因为电势是标量，
积分通常比直接分解电场容易；求出 $\varphi$ 后再用 $\vect E=-\nabla\varphi$ 得到电场。若场点
落在连续分布内部，要注意积分核在邻近电荷元处是否可积；对理想点电荷，自能部分会发散，不能
未经说明地混入连续分布能量。

均匀带电薄球壳是典型例子。球外电势等于把全部电荷集中在球心：
\[
 \varphi_{\rm out}(r)=\frac1{4\pi\varepsilon_0}\frac Qr.
\]
球内电场为零，所以球内电势处处等于表面值
\[
 \varphi_{\rm in}(r)=\frac1{4\pi\varepsilon_0}\frac QR.
\]
电势在球面连续，法向电场却因面电荷发生跃变。对于均匀带电实心球，利用 Gauss 定律先求球内
场，再从表面积分到内部，可得
\[
 \varphi(r)=\frac1{4\pi\varepsilon_0}\frac{Q}{2R}
 \left(3-\frac{r^2}{R^2}\right),\qquad r\le R.
\]
在 $r=R$ 处，它与球外电势及其一阶导数连续；在球心电势取得最大值但电场为零。

\subsection*{Poisson 方程与电势图像}

把 $\vect E=-\nabla\varphi$ 代入 Gauss 定律的微分形式，得到
\[
 \boxed{\nabla^2\varphi=-\frac{\rho}{\varepsilon_0}}.
\]
无电荷区域满足 Laplace 方程。电势曲面的弯曲程度由当地电荷密度决定，但电势本身还受整个空间
电荷和边界条件影响。一个区域内 $\rho=0$ 并不表示 $\varphi=0$ 或 $\vect E=0$；它只表示电势
在该区域为调和函数。

沿一条电场线，取弧长正方向与 $\vect E$ 同向，有
$\dd\varphi/\dd s=-E<0$，所以电势必沿场线降低。等势线画得越密，表示相同电势差发生在越短
距离内，场强越大；但只有在相邻等势线取相同电势间隔时，线距才能用于比较场强。图示若没有标明
等势值，线的疏密只能作定性提示，不能直接当成定量比例尺。
